\documentclass[11pt,a4paper]{article}

\usepackage[french]{babel}
\usepackage[utf8]{inputenc}
\usepackage[T1]{fontenc}
\usepackage{lmodern}
\usepackage[margin=2.5cm]{geometry}
\usepackage{amsmath}
\usepackage{booktabs}
\usepackage{longtable}
\usepackage{array}
\usepackage{xcolor}
\usepackage{hyperref}
\usepackage{fancyhdr}
\usepackage{enumitem}

\definecolor{solidateal}{RGB}{15,76,79}
\definecolor{gris}{RGB}{90,90,90}
\definecolor{vertaccord}{RGB}{45,110,80}
\definecolor{orangevigilance}{RGB}{180,120,20}
\definecolor{rougerefus}{RGB}{160,50,45}

\hypersetup{
  colorlinks=true,
  linkcolor=solidateal,
  urlcolor=solidateal,
  pdftitle={SOLIDA -- Explication du modèle de scoring},
  pdfauthor={SOLIDA}
}

\pagestyle{fancy}
\fancyhf{}
\fancyhead[L]{\small\textcolor{gris}{SOLIDA -- Explication du modèle de scoring}}
\fancyhead[R]{\small\textcolor{gris}{\thepage}}
\renewcommand{\headrulewidth}{0.4pt}

\newcommand{\brk}{\allowbreak}
\newcommand{\fichier}[1]{\noindent\textbf{Implémentation :} \tiny\ttfamily#1\normalfont\normalsize\par}
\newcolumntype{L}[1]{>{\raggedright\arraybackslash}p{#1}}

\title{\textcolor{solidateal}{\Huge SOLIDA} \\[0.3em] \Large Comment le modèle calcule un score, et ce que chaque variable y change}
\author{}
\date{Dernière mise à jour : \today}

\begin{document}

\maketitle

\section*{Objet de ce document}

Ce document explique, en détail et sans raccourci, comment le moteur de scoring de SOLIDA
transforme les données d'un sociétaire en un score, une décision et une fiche de justification.
Il répond à trois questions précises : comment le modèle considère-t-il les variables et leurs
poids, comment un point de score se calcule-t-il exactement, et quel est l'impact réel de chaque
variable sur la sortie du modèle. Chaque chiffre cité ici est extrait du code réel ou des
artefacts d'entraînement réels du projet -- aucune valeur n'est approximée ou inventée pour les
besoins de l'explication.

\tableofcontents

\section{Vue d'ensemble architecturale}

\subsection{Un modèle additif, pas une boîte noire}

Le modèle utilisé est un \textbf{EBM} (\emph{Explainable Boosting Machine}), une variante moderne
des modèles additifs généralisés (GAM). Contrairement à une régression logistique classique, il
n'a pas un coefficient unique par variable : chaque variable dispose de sa \textbf{propre courbe
de score} $f_j$, apprise par petites étapes de boosting successives, et le résultat final est la
somme de toutes ces courbes :

\begin{equation}
\text{log-odds} = \beta_0 + f_1(x_1) + f_2(x_2) + \dots + f_n(x_n)
\end{equation}

où $\beta_0$ est l'intercept du modèle et chaque $f_j$ une fonction non linéaire propre à la
variable $x_j$ (une table de valeurs par tranche pour une variable continue, une valeur par
modalité pour une variable catégorielle). Cette structure purement additive est ce qui rend le
modèle explicable par construction : il n'y a aucune interaction cachée à décoder, la contribution
de chaque variable à la décision finale est un terme séparé et directement lisible.

\subsection{Un choix architectural délibéré}

Le projet exclut explicitement toute approche qui réintroduirait de l'opacité : pas de
vectorisation de texte libre, pas de NLP, pas de réseau de neurones profond sur les variables
structurées. La documentation de cadrage du projet le formule ainsi :

\begin{quote}\itshape
« Pas de vectorisation ni de NLP : ce serait introduire une opacité que l'architecture EBM a
précisément été choisie pour éviter, sur un problème qui n'en a pas besoin. »
\end{quote}

Toutes les variables d'entrée sont donc des grandeurs structurées (montants, ratios, compteurs,
catégories fermées) -- jamais un champ de texte libre ni un document scanné analysé
automatiquement.

\section{Le pipeline complet, de la donnée brute au score affiché}

Cinq étapes successives, chacune vérifiable indépendamment.

\subsection{Étape 1 -- Construction des features}

Pour chaque demande, le backend calcule un dictionnaire de variables à partir de l'historique
CORE-SIM du sociétaire (épargne, crédits antérieurs, profil) et, le cas échéant, de son groupe de
caution. Ce calcul est strictement identique à celui utilisé pendant l'entraînement -- les mêmes
fonctions Python pures sont partagées entre le pipeline d'entraînement (\texttt{modelisation/}) et
le moteur d'inférence (\texttt{backend/}), pour garantir qu'aucun écart de définition ne puisse
exister entre ce que le modèle a appris et ce qu'il reçoit en production.

\subsection{Étape 2 -- Le modèle EBM produit un log-odds}

Le modèle évalue chaque courbe $f_j$ à la valeur reçue pour la variable correspondante, puis les
additionne avec l'intercept $\beta_0$ (valeur réelle du modèle SOCLE actuel : $\beta_0 = -3{,}074$).
Le résultat est un log-odds brut, pas encore une probabilité.

\subsection{Étape 3 -- Calibration}

Le log-odds brut est recalé par un calibrateur de type Platt (une régression logistique simple
appliquée sur la sortie du modèle), pour que la probabilité produite corresponde réellement aux
fréquences observées dans les données -- et pas seulement à un bon classement relatif des
dossiers. Cette étape est ce qui permet d'obtenir une erreur de calibration (ECE) mesurée à
$0{,}011$ pour le modèle SOCLE et $0{,}018$ pour le modèle enrichi : des valeurs basses, signe
d'une probabilité fiable en valeur absolue et pas seulement en classement.

\subsection{Étape 4 -- Transformation en score (méthode PDO)}

La probabilité calibrée est ensuite transformée en un score compris entre 300 et 850, selon la
méthode \textbf{PDO} (\emph{Points to Double the Odds}) -- la méthode standard de l'industrie du
crédit, la même famille de méthode qui produit les scores FICO :

\begin{align}
\text{log-odds}_{\text{bon}} &= \ln\!\left(\frac{1-p}{p}\right) \\
\text{facteur} &= \frac{\text{PDO}}{\ln 2} \\
\text{décalage} &= \text{score}_{\text{réf}} - \text{facteur} \times \ln(\text{odds}_{\text{réf}}) \\
\text{score} &= \text{décalage} + \text{facteur} \times \text{log-odds}_{\text{bon}}
\end{align}

avec $p$ la probabilité de défaut (convention bancaire : score élevé = risque \emph{faible}, d'où
le rapport $(1-p)/p$ et non son inverse). Les paramètres réellement configurés dans SOLIDA
aujourd'hui sont :

\begin{center}
\begin{tabular}{ll}
\toprule
Paramètre & Valeur actuelle \\
\midrule
PDO (points pour doubler les chances) & 20 \\
Score de référence & 600 \\
Rapport de cotes de référence & 50 \\
Score minimum / maximum & 300 / 850 \\
\bottomrule
\end{tabular}
\end{center}

Concrètement : \textbf{chaque fois que le rapport (chances de bien rembourser) / (chances de faire
défaut) double, le score augmente de 20 points}. Le score final est borné à $[300, 850]$ : un
dossier extrême ne peut jamais sortir de cette plage, même si le calcul brut le dépasserait.

\subsection{Étape 5 -- Décomposition en points par variable}

C'est l'étape qui produit ce qu'affiche la fiche de justification. La transformation ci-dessus
étant \textbf{affine} (une mise à l'échelle plus un décalage), chaque terme $f_j(x_j)$ du modèle
se convertit indépendamment en points, par le même facteur d'échelle :

\begin{align}
\text{points de base} &= \text{décalage} + \text{facteur} \times \beta_0 \\
\text{points}_j &= \text{facteur} \times f_j(x_j) \quad \text{pour chaque variable } j
\end{align}

Le score final affiché est garanti égal à la somme de tous ces termes :

\begin{equation}
\text{score} = \text{points de base} + \sum_j \text{points}_j
\end{equation}

Cette égalité n'est pas une simple convention d'affichage : elle est \textbf{vérifiée
automatiquement à chaque calcul}, et le système \emph{rejette} le scoring si la somme ne
correspond pas exactement au score annoncé (tolérance de 1 point, pour absorber les seuls
arrondis flottants). Une fiche de justification dont les points ne somment pas au score détruirait
la crédibilité du produit devant un agent ou un auditeur -- le système préfère refuser un calcul
incohérent plutôt que de le montrer.

\fichier{backend/solida/domain/rules/scorecard.py}

\section{De la probabilité à la décision : la grille}

Une fois la probabilité de défaut $p$ calculée, la tranche de décision (accord, accord sous
condition, comité de crédit, refus) se détermine par comparaison à un \textbf{seuil économique},
lui-même issu de la théorie de la classification sensible aux coûts (Elkan, 2001) :

\begin{equation}
\text{seuil} = \frac{\text{marge}}{\text{marge} + \text{LGD}}
\end{equation}

où LGD (\emph{Loss Given Default}) est la perte proportionnelle en cas de défaut, et marge la
marge économique sur un crédit remboursé normalement. Avec les valeurs actuellement configurées
(marge $= 0{,}15$, LGD $= 0{,}75$), le seuil économique vaut $0{,}15 / 0{,}90 \approx 0{,}167$. Ce
seuil est ensuite modulé par des multiplicateurs de zone pour définir quatre bandes :

\begin{center}
\begin{tabular}{lll}
\toprule
Condition sur $p$ & Tranche & Multiplicateur actuel \\
\midrule
$p < \text{seuil} \times 0{,}6$ & \textcolor{vertaccord}{\textbf{Accord}} & 0,6 \\
$p < \text{seuil} \times 1{,}0$ & \textcolor{orangevigilance}{\textbf{Accord sous condition}} & 1,0 \\
$p < \text{seuil} \times 1{,}6$ & \textbf{Comité de crédit} & 1,6 \\
sinon & \textcolor{rougerefus}{\textbf{Refus}} & -- \\
\bottomrule
\end{tabular}
\end{center}

Ces seuils et multiplicateurs sont \textbf{modifiables sans redéploiement}, par la supervision, via
l'écran de paramétrage de la grille -- ce ne sont jamais des constantes figées dans le code.

\fichier{backend/solida/domain/rules/grille.py}

\section{Catalogue complet des variables}

Deux catalogues coexistent. Le \textbf{catalogue SOCLE} (20 variables) s'applique à tout dossier.
Le \textbf{catalogue enrichi} (SOCLE + 6 variables de groupe) s'applique uniquement quand le
sociétaire appartient à un groupe de caution solidaire jugé éligible par la cascade de décision --
un mécanisme séparé, non traité ici.

\subsection{Catalogue SOCLE (20 variables)}

\begin{longtable}{L{6cm}L{9cm}}
\toprule
\textbf{Variable} & \textbf{Nature} \\
\midrule
\endhead
anciennete\_societaire\_mois & Ancienneté du sociétaire dans la coopérative, en mois \\
nb\_personnes\_a\_charge & Nombre de personnes à charge déclarées \\
zone\_residence & Zone de résidence (urbain / rural) \\
parts\_sociales\_montant & Montant des parts sociales détenues \\
revenu\_mensuel\_declare & Revenu mensuel déclaré \\
ratio\_endettement & Mensualité du crédit demandé rapportée au revenu \\
duree\_demandee\_mois & Durée demandée pour le crédit \\
solde\_epargne\_moyen\_6m & Solde d'épargne moyen sur les 6 derniers mois \\
nb\_mois\_avec\_depot\_12m & Nombre de mois avec au moins un dépôt sur 12 mois \\
tendance\_epargne\_12m & Tendance de l'épargne sur 12 mois (hausse / stable / érosion) \\
volatilite\_epargne & Volatilité du solde d'épargne \\
ratio\_epargne\_revenu & Épargne moyenne rapportée au revenu \\
anciennete\_epargne\_mois & Ancienneté de la relation d'épargne \\
ratio\_epargne\_montant & Épargne moyenne rapportée au montant demandé \\
nb\_credits\_anterieurs & Nombre de crédits antérieurs \\
nb\_incidents\_anterieurs & Nombre d'incidents de remboursement antérieurs \\
max\_jours\_retard\_historique & Retard maximal observé sur l'historique \\
montant\_max\_rembourse & Plus gros montant déjà remboursé avec succès \\
numero\_cycle & Numéro du cycle de crédit en cours \\
ratio\_montant\_historique & Montant demandé rapporté au meilleur remboursement passé \\
\bottomrule
\end{longtable}

\subsection{Variables additionnelles du catalogue enrichi (+6)}

\begin{longtable}{L{6cm}L{9cm}}
\toprule
\textbf{Variable} & \textbf{Nature} \\
\midrule
\endhead
taille\_groupe & Nombre de membres actifs dans le groupe de caution \\
anciennete\_groupe\_mois & Ancienneté du groupe de caution \\
nb\_credits\_groupe\_anterieurs & Nombre de crédits antérieurs portés par le groupe \\
nb\_incidents\_groupe\_anterieurs & Nombre d'incidents de remboursement du groupe \\
max\_jours\_retard\_groupe\_6m & Retard maximal du groupe sur les 6 derniers mois \\
nb\_cautions\_appelees\_anterieures & Nombre de cautions déjà appelées dans le groupe \\
\bottomrule
\end{longtable}

\fichier{modelisation/data/modeles/bundle\_socle/manifeste.json, bundle\_enrichi/manifeste.json}

\section{Impact réel de chaque variable}

\subsection{Importance globale, calculée sur les données réelles d'entraînement}

Le tableau ci-dessous n'est pas une estimation : il est calculé directement à partir du modèle
SOCLE entraîné et des 12\,690 dossiers réels du jeu d'entraînement (\texttt{modelisation/\brk
data/\brk processed/\brk jeu\_socle.parquet}), via l'explication globale native d'EBM (moyenne de
la valeur absolue de la contribution de chaque variable sur l'ensemble du jeu de données).

\begin{center}
\begin{tabular}{lr}
\toprule
\textbf{Variable} & \textbf{Poids relatif global} \\
\midrule
Épargne rapportée au montant demandé & 0,73 \\
Taux d'endettement & 0,52 \\
Durée demandée & 0,33 \\
Revenu mensuel déclaré & 0,29 \\
Montant demandé / meilleur remboursement & 0,26 \\
Volatilité de l'épargne & 0,22 \\
Plus gros montant déjà remboursé & 0,15 \\
Épargne rapportée au revenu & 0,13 \\
Personnes à charge & 0,08 \\
Zone de résidence & 0,08 \\
Solde d'épargne moyen (6 mois) & 0,07 \\
Montant des parts sociales & 0,07 \\
Ancienneté sociétaire & 0,06 \\
Ancienneté de l'épargne & 0,06 \\
Régularité des dépôts (12 mois) & 0,04 \\
\bottomrule
\end{tabular}
\end{center}

Ce classement révèle que le modèle fonde l'essentiel de sa décision sur deux familles de signaux
: la \textbf{capacité de remboursement} (taux d'endettement, durée, revenu) et la \textbf{capacité
d'épargne relative à la demande} (épargne rapportée au montant, à l'historique de remboursement).
Les variables purement démographiques (personnes à charge, zone de résidence) pèsent nettement
moins.

\subsection{Importance globale et contribution individuelle : deux choses différentes}

Ce classement dit ce qui compte \emph{en moyenne} sur tout le portefeuille. Il ne prédit pas ce
qui va compter pour un dossier précis : une variable secondaire dans ce classement peut peser
lourd sur un dossier particulier si sa valeur y est extrême (une volatilité d'épargne
anormalement haute, par exemple, peut faire perdre beaucoup de points à un dossier donné même si
cette variable pèse peu en moyenne). C'est précisément pour cette raison que la fiche de
justification affiche toujours la contribution \textbf{réelle et individuelle} de chaque variable
pour le dossier en cours, jamais un classement générique plaqué sur tous les dossiers.

\section{Couverture par rapport aux 5C du crédit}

Le cadre classique d'analyse du risque de crédit distingue cinq dimensions (les « 5C »). Voici ce
que le modèle couvre réellement aujourd'hui, variable par variable :

\begin{center}
\begin{tabular}{L{2.7cm}L{1.8cm}L{9cm}}
\toprule
\textbf{Dimension} & \textbf{Couvert.} & \textbf{Variables concernées} \\
\midrule
Capacity (capacité de remboursement) & Forte & revenu\_mensuel\_declare, ratio\_endettement, duree\_demandee\_mois, ratio\_montant\_historique, montant\_max\_rembourse \\
Capital (fonds propres) & Forte & solde\_epargne\_moyen\_6m, parts\_sociales\_montant, ratio\_epargne\_revenu, anciennete\_epargne\_mois, volatilite\_epargne \\
Character (comportement, réputation) & Forte & nb\_credits\_anterieurs, nb\_incidents\_anterieurs, max\_jours\_retard\_historique, numero\_cycle, anciennete\_societaire\_mois, nb\_mois\_avec\_depot\_12m, et en enrichi le comportement du groupe \\
Collateral (garantie) & Partielle & ratio\_epargne\_montant (épargne nantie) et, en mode enrichi, la caution solidaire du groupe -- aucune garantie matérielle classique n'est modélisée \\
Conditions (contexte économique) & Faible & zone\_residence uniquement ; l'objet du crédit est collecté mais n'est pas une variable du modèle, et aucune donnée macro-économique ou sectorielle n'entre dans le calcul \\
\bottomrule
\end{tabular}
\end{center}

\section{Limites et garde-fous encore ouverts}

Par souci d'honnêteté, cette section documente ce que le modèle \textbf{ne garantit pas encore},
sur la base des artefacts d'audit réellement produits par le pipeline d'entraînement lui-même --
pas d'une opinion externe.

\subsection{Données d'entraînement synthétiques}

Le modèle est entraîné sur 12\,690 dossiers générés par un simulateur
(\texttt{simulateur/}), pas sur l'historique réel de la coopérative. C'est documenté et assumé
dans tout le projet : c'est un prototype conçu pour être recalibré sur données réelles avant tout
usage en production.

\subsection{Un seuil de performance qui déclenche sa propre alerte}

Le fichier d'audit produit automatiquement à l'entraînement (\texttt{modelisation/\brk data/\brk
modeles/\brk socle\_audit\_fuite.json}) contient :

\begin{quote}\ttfamily\footnotesize
"auc\_test": 0.898,\\
"audit\_renforce\_requis": true,\\
"motif": "AUC test strictement supérieure\\
à 0,88 : revue humaine de fuite requise."
\end{quote}

Les contrôles automatiques (absence de colonne interdite, pas de fuite temporelle, classes cible
correctes) sont passés. Mais la \textbf{revue humaine que ce seuil déclenche explicitement n'a
laissé aucune trace de validation} dans le dépôt à ce jour -- c'est une étape identifiée comme
nécessaire, pas encore réalisée.

\subsection{Équité et seuil de sélection non actés}

Le fichier \texttt{socle\_revue\_equite.json} signale de la même façon deux points en attente :

\begin{quote}\ttfamily\footnotesize
"seuil\_selection\_non\_defini": true,\\
"revue\_humaine\_ecart\_\brk risque\_superieur\_5\_points": true
\end{quote}

Le seuil de décision opérationnel n'est pas figé, et un écart de risque entre segments dépasse 5
points sans qu'une revue d'équité n'ait été actée à ce jour.

\subsection{Ce que cela signifie concrètement}

Le modèle est \textbf{techniquement solide} : architecture explicable par construction, bonne
discrimination, bonne calibration mesurée (Brier $0{,}046$, ECE $0{,}011$--$0{,}018$). Mais deux
garde-fous que le projet s'est lui-même explicitement imposés restent ouverts, et les données
d'entraînement ne reflètent pas encore la réalité de la coopérative. Présenter ce modèle comme
définitivement validé serait inexact ; le présenter comme un prototype rigoureux et transparent,
avec une feuille de route claire vers la validation, est la description honnête de son état
actuel.

\end{document}
